\begin{question}

A beam of spin $\frac 1 2$ atoms goes through a series of
Stern–Gerlach type measurements as follows.

\begin{enumerate}[label=\alph*.]

\item The first measurement accepts $S_z = \frac{\hbar}{2}$ atoms and
  rejects $S_z = -\frac{\hbar}{2}$ atoms.

\item The second measurement accepts $S_n = \frac{\hbar}{2}$ atoms and
  rejects $S_n = -\frac{\hbar}{2}$, where $S_n$ is the eigenvalue of
  the operator $\vectorbold{S} \cdot \vectorbold{\hat{n}}$, with
  $\vectorbold{\hat{n}}$ making an angle $\beta$ in the xz-plane with
  respect to the z-axis.

\item The third measurement accepts $S_z = \frac{\hbar}{2}$ atoms and
  rejects $S_z = -\frac{\hbar}{2}$ atoms.

\end{enumerate}

\begin{questionpart}
  What is the intensity of the final $S_z = -\frac{\hbar}{2}$ beam
  when the $S_z = \frac{\hbar}{2}$ beam surviving the first
  measurement is normalized to unity?
\end{questionpart}

\begin{questionpart}
 How must we orient the second measuring apparatus if we are to
 maximize the intensity of the final $S_z = \frac{\hbar}{2}$ beam?
\end{questionpart}

\end{question}

\begin{purpose}
  This one looks like basic practice with how state collapses across
  observations and how that affects probabilities.
\end{purpose}

\begin{answer}

  \begin{answerpart}
    For simplicity, let's put up a quick diagram of the pipeline of
    experiments:

    \begin{figure}
      \inlabel{pipeline}
      \includegraphics[width=3in]{pipeline.png}
      \caption{Phase a}
    \end{figure}

  \end{answerpart}

  \begin{answerpart}
    bar
  \end{answerpart}

\end{answer}
